\documentclass[]{scrartcl}
\usepackage{graphicx}
\usepackage{color}
\usepackage[ngerman]{babel}
\usepackage{hyperref}
\usepackage{calc} 
\usepackage{enumitem}
\usepackage{titlesec}
%\pagestyle{headings}
%\pagestyle{headings}

\begin{document}

\title{
	\includegraphics*[width=0.75\textwidth]{images/hu_logo.png}\\
	\vspace{24pt}
	Einf"uhrung in die \\ Architektur der Moderne\\ Konzeptionen, Imaginationen,\\ mediale Rezeption}
\subtitle{Proseminar WS 16/17\\
          Prof. Dr. Kai Kappel\\
          Institut f"ur Kunst- und Bildgeschichte \\ 
          Humboldt Universit"at zu Berlin}
\author{Lennard Wolf\\
        \href{mailto:lennard.wolf@student.hu-berlin.de}{lennard.wolf@student.hu-berlin.de}}
\maketitle
\begin{abstract}

Das einf"uhrende Seminar er"offnet kritische Zug"ange zur Begriffsbildung, Theorie und den wichtigsten Str"omungen der modernen Architektur. 

\end{abstract}
\newpage

\tableofcontents

\listoffigures
\newpage


\section{Einf"uhrende Sitzung\\(18.10.16)}
\subsection{Organisatorisches}

\begin{itemize}
    \item Erste beide Semester sind Grundlagenmodule und zu belegen
    \item Klausur wird geschrieben, erste und einzige in 1. Sem
    \item Struktur: "Uberblicksvorlesung + Proseminar + Tutorium
    \item MAP: Referat mit Team (ganzer Monat; MAP wird bei Agnes angemeldet)
    \item Moodle pw: \emph{Stahltr"ager}
    \item Pflichtsprechstunde vor dem Vortrag
\end{itemize}

Vortrag: \emph{Le Corbusiers Konzept einer Ville Contemporaine}; 20min, 1 Handout A4
Aufgabe dabei: Sagen, was die Bauten verk"orpern, die Ideengeschichte dahinter. Ob erst geb"aude und dann idee oder andersherum h"angt vom Vortrag ab.
Auch Kritik an m"oglichen Erkl"arungen aus Geschichtsb"uchern "uben (wenn angebracht).

\subsection{Einf"uhrung}
\subsubsection{Definitionen der Moderne?}
\begin{itemize}
    \item Haeufig sagt man: nach dem ersten Weltkrieg 1918 (Bauhaus)
    \item Aber: Jugendstil, Reformarchitektur, Eisen \& Stahlkonstruktionen etc..?
    \item Baustoffe, Schmuck etc.
    \item Seminar setzt im 19.Jh. an
\end{itemize}

\subsubsection{Wann endet die Moderne?}

\begin{itemize}
    \item Postmoderne, Zweite Moderne etc.?
    \item \emph{Modern is not contemporary}
\end{itemize}

\section{Wie beschreibe und analysiere ich moderne Architektur?\\(25.10.16)}

Treffen an der Ged"achtniskirche


\section{Zum Modernebegriff (nicht nur) in der Architektur\\(01.11.16)}

\subsection{Vornotizen}
\textbf{Definitionen und Abgrenzungen von Moderne}

\begin{description}[leftmargin=!,labelwidth=\widthof{\bfseries P2}]
  \item[Greenberg (Arch.)] \emph{Moderne Architektur bedeutet -- vereinacht gesagt -- funktionelle, geometrische Strenge und das Vermeiden von Dekoration und Ornament.} | Kap. `Modern und postmodern', Die Essenz der Moderne 
  \item[Greenberg (Kunst)] Bezeichnet K"unstler der Moderne als \emph{unbeugsamen Helden, der sich auf keinerlei Kompromisse mit dem Bestehenden einl"a\ss t und allen Verlockungen der eleganten Welt edelm"utig widersteht}; \emph{Bedeutende Werke wird ein K"unstler nur zu Wege bringen, wenn er alle au\ss erk"unstlcihen Anspr"uche und Erwartungen abweist, um seine Arbeit ganz aus sich sekbst heraus zu begr"unden.} $\rightarrow$ Das Publikum wurde dem Autor nach immer verst"andnisloser. | \emph{..selbstkritische Besinnung der Kunst auf ihre ureigenen und spezifischen M"oglichkeiten ist, Greenberg zufolge, das Charakterisikum des `Modernismus' }
  
  \item[Kretschmer] Latein: \emph{der Gegenwart angeh"orig} | im 19.Jh. als Gegensatz zu \emph{vergangen} | `Siegeszug' erst nach WWII | \emph{Die Moderne ist generell eine von technischem Positivismus gekennzeichnete Emanzipationsbewegung zur Befreiung aus der Fremdbestimmung.} | \emph{Es ging um eine m"oglichst gute Versorgung der neuen Massengesellschaft.} 
\end{description}

\textbf{Besprechung der Notizen}

\begin{itemize}
  \item Probleme zum einen der Datierung und der begrifflichen Festsetzung
  \item Es ist schwer w"ahrend einer Zeit "uber diese Zeit selber zu reden
  \item Freigang: Es gibt keine gro\ss e Erz"ahlung! Es besteht aus vielerlei Nebenwegen.
  \item Es ist kein Stil, kein formaler Rahmen, keine feste Zeit
\end{itemize}


\subsection{Sitzungsvortrag}

\begin{itemize}
  \item Seit Sp"atantike (5.Jh.): Die noch selbst erlebt Zeit
  \item Neuzeit: Kunst/Architektur die nicht dezidiert von der antiken Bautradition ableitbar ist (Schon Hauch von Avantgarde)
  \item Moderne in Geisteswissenschaften: industrielle Revolution, AUfkl"arung und S"akularisierung | Moderne in Philosophie: Aufkl"arung
  \item Beginn zwischen sp"aten 18. und mittleren 19. Jh. $\rightarrow$ Zeit des "Ubergangs von einem feudalistischen zu einem b"urgerlichen Gesellschaftsmodell
  \item Erste Erw"ahung 1886 von Eugen Wolff im Kontext der Literatur
  \item Moderne ist kein Stil- oder Epochenbegriff, sondern die Bezeichnung einer intellektuell wie formal (materiell wie gestalterisch) au\ss erordentlich vielf"altigen Str"omung.
  \item Bedingt durch die z.T. neuen, z.T. nun industriell herstellbaren Materialien Eisen/Stahl, Glas und sp"ater Beton: Industiebauten, Verkehrsbauten, Messebauten des 19. und 20. Jh.s (Paxton, Fr\`{e}res, Eiffel: \emph{Historische Moderne}) 
  \item Hochhausbauten des sp"aten 19. und erszen Drittel des 20. Jh.s in den USA; Arts and Crafts Architektur; Reformarchitektur (Jugendstil bzw. Wiener Sezession, Gartenstadtbewegung, Deutscher Werkbund, Darmst"adter Mathildenh"ohe, heimatbezofene Architektur)
  \item Klassiche Moderne 1910-1933 (Avantgardistische Architektur, Expressionistische, Organische Architektur, Traditionalismus) und viele mehr...
  \item Metabolismus, Cities:Moving
\end{itemize}


\section{Stahl und Glaskonstruktionen: Weltausstellungen des 19. Jahrhunderts\\(08.11.16)}

\section{Stahl und Glaskonstruktionen: Hochhausbauten\\(15.11.16)}

\section{Buildings to live in, not to look at. Absagen an die Gliederungen des Historismus\\(22.11.16)}

Bei Postmoderne: McMansions!


\newpage
\section{"Uber den Professor}
Prof. Dr. Kai Kappel
%\begin{figure}[h]
%	\centering
%	\includegraphics[width=0.5\textwidth]{images/template.png}
%	\caption{Template Bild}
%	\label{fig:template}
%\end{figure}


\end{document}
